\section{Sum of Beauty of All Substrings}
\label{sec:str-sum-of-beauty}

\problemstatement{The beauty of a string is the difference between the
highest and lowest character frequencies in it (considering only characters
that appear). Compute the sum of the beauty of all substrings.}

\begin{patternbox}
There are $O(n^2)$ substrings, so an $O(n^2)$ approach is expected: fix each
start, then \textbf{extend the end one character at a time, maintaining a
frequency array}. After each extension, the beauty (max minus min nonzero
frequency) is computed and added.
\end{patternbox}

\subsection*{Grow substrings, update frequencies incrementally}

For each starting index $i$, reset a frequency array and extend the end $j$
from $i$ to the end. Adding $s[j]$ updates one frequency in $O(1)$. The
beauty of substring $[i,j]$ is the maximum frequency minus the minimum
\emph{nonzero} frequency. Scanning the 26 counts to find those extremes is
$O(26)$ per substring, giving $O(26n^2)$ overall -- fast enough for the
constraints.

\begin{ideabox}[Reuse the Frequency Array as the Window Grows]
Fixing the start and extending the end lets each new substring's frequencies
be obtained by a single increment, avoiding a recount. Only the max/min scan
over the fixed 26-letter alphabet is repeated, keeping the inner cost
constant.
\end{ideabox}

\subsection*{Algorithm}

\begin{enumerate}
  \item For each start $i$: clear a $26$-count array.
  \item For each end $j\ge i$: increment \texttt{count[s[j]]}; compute
        beauty $=\max - \min_{\text{nonzero}}$ over the counts; add to the
        total.
  \item Output the total.
\end{enumerate}

\subsection*{Dry run}

\dryrun{\texttt{"aabcb"}.}

\begin{itemize}
  \item Most substrings have beauty $0$ (uniform or single char);
        \texttt{"aabcb"} and \texttt{"aab"} contribute positive beauty.
  \item Summing all substrings' beauties gives $5$.
\end{itemize}

\subsection*{C++ implementation}

\begin{lstlisting}
#include <bits/stdc++.h>
using namespace std;

int beautySum(const string& s) {
    int n = s.size(), total = 0;
    for (int i = 0; i < n; i++) {
        int cnt[26] = {0};
        for (int j = i; j < n; j++) {
            cnt[s[j] - 'a']++;
            int mx = 0, mn = INT_MAX;
            for (int k = 0; k < 26; k++) if (cnt[k]) {
                mx = max(mx, cnt[k]);
                mn = min(mn, cnt[k]);
            }
            total += mx - mn;
        }
    }
    return total;
}

int main() {
    cout << beautySum("aabcb") << "\n";            // 5
    return 0;
}
\end{lstlisting}

\begin{complexitybox}
\textbf{Time Complexity.} $O(26\,n^2)$ -- all substrings, each an $O(26)$
extreme scan. \textbf{Space Complexity.} $O(1)$ (a fixed count array).
\end{complexitybox}

\begin{takeawaybox}
When all $O(n^2)$ substrings must be examined, fix the start and extend the
end while maintaining an incremental frequency array -- only the small
fixed-alphabet extreme scan repeats. This ``grow the window, update counts''
pattern recurs in substring-statistic problems.
\end{takeawaybox}
